% ============================================================
%  系统开发工具基础 · 实验报告 
%  作者：曾奕豪
%  编译方式：XeLaTeX（需 ctex 宏包支持中文）
%  用法：每周复制本文件为 weekN.tex，仅修改周次、日期、内容即可
%  说明：本模板不依赖任何图片即可编译；
%        图片文件存在时自动插入，不存在时显示占位提示框。
%  图片统一放在 paper/figures/weekN_picture/ 子文件夹中（N 为周次），
%        命名规则 weekN_xxx.png；每周只需修改下面的 \Week 即可。
% ============================================================

\documentclass[12pt,a4paper]{article}

% ---------- 中文字体与页面 ----------
% 中文支持：使用 ctex 宏包（编译用 XeLaTeX）。
% 若在 Linux/TeX Live 下缺少 ctex，可安装 texlive-lang-chinese。
\usepackage[UTF8]{ctex}
\usepackage[margin=2.5cm]{geometry}
\usepackage{graphicx}
\usepackage{float}
\usepackage{booktabs}
\usepackage{xcolor}
\usepackage{listings}
\usepackage{enumitem}
\usepackage[colorlinks=true,linkcolor=black,urlcolor=blue]{hyperref}

% ---------- 页眉 / 页脚 ----------
\usepackage{fancyhdr}
\pagestyle{fancy}
\fancyhf{}
\fancyhead[R]{实验报告 · 第 \Week 周}
\fancyfoot[C]{\thepage}
\renewcommand{\headrulewidth}{0.4pt}

% ---------- 代码高亮环境 ----------
\lstset{
  basicstyle=\ttfamily\small,
  keywordstyle=\color{blue!70!black}\bfseries,
  commentstyle=\color{gray!70!black}\itshape,
  stringstyle=\color{red!60!black},
  numbers=left,
  numberstyle=\tiny\color{gray},
  frame=single,
  breaklines=true,
  tabsize=4,
  captionpos=b,
  showstringspaces=false,
}

% ---------- 图片插入宏 ----------
% 题内图片：图注手动指定文字，不自动编号。
% 图片统一放在 paper/figures/week\Week_picture/ 子文件夹中，
% 文件名建议带周次前缀（如 week1_q1.png），每周只需修改 \Week 即可。
%   \resetpic              每题开头调用一次（保留自 week1，兼容旧结构）
%   \resultimg[宽度]{图片文件名}{图注}   图片存在→插入；不存在→占位框（不报错）
% 自定义图注（如 commit 截图）：
%   \resultimgcap[宽度]{图片文件名}{图注}
\newcounter{picnum}
\newcommand{\resetpic}{\setcounter{picnum}{0}}
% 图片目录：根据 \Week 自动定位到 weekN_picture 子文件夹
\newcommand{\picdir}{../../figures/week\Week_picture/}
\newcommand{\resultimg}[3][0.8\textwidth]{%
  \IfFileExists{\picdir #2}{%
    \begin{figure}[H]
      \centering
      \includegraphics[width=#1]{\picdir #2}
      \parbox{0.8\textwidth}{\centering\small #3}
    \end{figure}%
  }{%
    \begin{center}
      \fbox{\parbox{0.8\textwidth}{\centering\itshape [图片待补充] \detokenize{#2}\\\small #3}}
    \end{center}%
  }%
}
\newcommand{\resultimgcap}[3][0.8\textwidth]{%
  \IfFileExists{\picdir #2}{%
    \begin{figure}[H]
      \centering
      \includegraphics[width=#1]{\picdir #2}
      \parbox{0.8\textwidth}{\centering\small #3}
    \end{figure}%
  }{%
    \begin{center}
      \fbox{\parbox{0.8\textwidth}{\centering\itshape [图片待补充] \detokenize{#2}\\\small #3}}
    \end{center}%
  }%
}

% ---------- 文档信息（每周只需改这三处） ----------
\newcommand{\Week}{4}                        % 周次
\newcommand{\ReportTitle}{系统开发工具基础实验报告}
\newcommand{\ReportDate}{\today} % 日期，编译时自动取当天
\newcommand{\RepoLink}{https://github.com/hzydy00/Base-of-system-development-tools} % 仓库链接

% ============================================================
\begin{document}

% ---------- 封面 / 标题区 ----------
\begin{center}
  {\LARGE\bfseries \ReportTitle}\\[6pt]
  {\Large 第 \Week\ 周}\\[12pt]
  {\large 姓名：曾奕豪}\\
  {\large 课程：系统开发工具基础}\\[6pt]
  {\large \ReportDate}\\[16pt]
  {\large 仓库：\url{\RepoLink}}
\end{center}
\vspace{0.8em}
\hrule
\vspace{1.2em}

% ---------- 1. 课上内容与结果 ----------
\section{课上内容与结果}
% 开头先简述本周实验的主题、目标与涉及的工具 / 技术，再逐题列出课上内容。
% 课上完成 4 个不同主题的题目，每题按「题目要求 → 代码 → 过程截图 → 实验结果」组织。
\subsection*{学习主题}

代码质量；元编程；大杂烩；

\subsection*{实验环境}
% 简述本实验的操作系统、命令行环境、所用工具及其版本。
本实验在 Windows 操作系统上完成，命令行环境使用 Git Bash（MINGW64）与 PowerShell。
Python 版本 3.14，代码质量与构建工具链为：ruff 0.16.6（代码格式化与静态检查）、
pytest 9.0.3（单元测试）、GNU Make 4.4.1（增量构建）、jq（JSON 处理）、curl（HTTP 请求）、Git 2.53.0.windows.2（版本控制）。

\subsection*{题目一}
\resetpic
% 题目一 问题截图（放在代码汇总之前）
\resultimgcap[0.8\textwidth]{week\Week_q1.png}{题目一 题目要求}

\begin{lstlisting}
## TODO: 题目一代码汇总，分为toml，python，bash三部分

# pyproject.toml 追加段
[tool.ruff]
line-length = 88
target-version = "py314"

[tool.pytest.ini_options]
testpaths = ["tests"]
pythonpath = ["src"]

\end{lstlisting}

\begin{lstlisting}[language=python]    
# tests/test_normal_name.py（新增）
import sys

from greetlab.cli import main


def test_normal_name_prints_greeting(capsys, monkeypatch):
    monkeypatch.setattr(sys, "argv", ["sdt-greet", "--name", "Alice"])
    main()
    captured = capsys.readouterr()
    assert captured.out == "Hello, Alice!\n"

\end{lstlisting}

\begin{lstlisting}[language=bash]    
# check.sh（新建）
#!/usr/bin/env bash
set -euo pipefail

ruff format --check .
ruff check .
python -m pytest


\end{lstlisting}

\begin{center}
{\large\bfseries 部分实验过程截图}
\end{center}

\resultimg{week\Week_check1.png}{题目一 实验过程截图}
\resultimg{week\Week_check1_2.png}{题目一 实验过程截图}

\begin{center}
{\large\bfseries 实验结果}
\end{center}

- ./check.sh 一次通过：ruff format --check（5 files already formatted）

-  ruff check（All checks passed!）

- pytest（2 passed）→ 退出码 0。

质量门禁绿灯，本地检查与提交前自检闭环成立。

\subsection*{题目二}
\resetpic
% 题目二 问题截图（放在代码汇总之前）
\resultimgcap[0.8\textwidth]{week\Week_q2.png}{题目二 题目要求}

\begin{lstlisting}
# TODO: 题目二代码汇总，分makefile和bash两部分

# makefile部分
.PHONY: all clean

all: stats.txt report.txt

stats.txt: data.csv stats.py
	python stats.py

report.txt: report.md stats.txt build_report.py
	python build_report.py

clean:
	rm -f stats.txt report.txt


\end{lstlisting}
\begin{lstlisting}
# python部分
# 首次构建 —— 预期依次执行两个配方
make
# 查看产物（预期 stats.txt=10；report.txt 含 "# Data Report" 和 "Total: 10"）
cat stats.txt
cat report.txt

# 无改动再 make —— 预期 "Nothing to be done for 'all'."
make

# touch 数据文件后再 make —— 预期两个配方重新执行
touch data.csv
make

# 清理 —— 预期只删除两个生成物
make clean
ls

\end{lstlisting}

\begin{center}
{\large\bfseries 部分实验过程截图}
\end{center}

\resultimg{week\Week_check2.png}{题目二 实验过程截图：首次构建，查看产物}
\resultimg{week\Week_check2_2.png}{题目二 实验过程截图：touch后再构建，清理产物}

\begin{center}
{\large\bfseries 实验结果}
\end{center}

- 首次 make：生成 stats.txt（10）和 report.txt（Total: 10）。

- 无改动再 make：Nothing to be done，配方不执行。

- touch data.csv 后 make：沿依赖链连锁重建两个产物。

- make clean：只删除生成物，源文件保留。
\subsection*{题目三}
\resetpic
% 题目三 问题截图（放在代码汇总之前）
\resultimgcap[0.8\textwidth]{week\Week_q3.png}{题目三 题目要求}

\begin{lstlisting}[language=bash]
# TODO: 题目三代码汇总
# api_report.sh
#!/usr/bin/env bash
set -euo pipefail

{
  echo "# Package Report"
  echo ""
  echo "| name | version | downloads |"
  echo "|---|---|---|"
  curl -fsS http://127.0.0.1:8000/packages.json \
    | jq -r '[.[] | select(.status == "active" and .downloads >= 100)]
             | sort_by(-.downloads, .name)
             | .[] | "| \(.name) | \(.version) | \(.downloads) |"'
} > summary.md

# 终端 A：启动本地服务（保持运行）
python -m http.server 8000

# 终端 B（另开一个）：取数 + 筛选排序，一步验证
curl -fsS http://127.0.0.1:8000/packages.json
curl -fsS http://127.0.0.1:8000/packages.json | jq -r '[.[] | select(.status == "active" and .downloads >= 100)] | sort_by(-.downloads, .name) | .[] | [.name, .version, .downloads] | @tsv'

# 运行脚本生成报告
./api_report.sh

# 查看结果
cat summary.md

# 回到终端 A，Ctrl+C 停止服务


\end{lstlisting}

\begin{center}
{\large\bfseries 部分实验过程截图}
\end{center}

\resultimg{week\Week_check3.png}{题目三 实验过程截图}

\begin{center}
{\large\bfseries 实验结果}
\end{center}

- python -m http.server 8000 启动后，curl -fsS 成功取回 packages.json  5 条记录。

- jq 筛选：beta（status=inactive） 与 epsilon（downloads=80 < 100） 被剔除，仅剩 lpha、gamma、delta 三条。

- 排序：downloads 降序为 450、450、120；gamma 与 delta 同为 450 时按 name 升序，delta 排在 gamma 前。

- api\_report.sh 生成 summary.md

\subsection*{题目四}
\resetpic
% 题目四 问题截图（放在代码汇总之前）
\resultimgcap[0.8\textwidth]{week\Week_q4.png}{题目四 题目要求}

\begin{lstlisting}
# TODO: 题目四代码汇总，分为2部分

# .gitignore
__pycache__/
*.pyc
.pytest_cache/
.ruff_cache/
*.egg-info/
dist/
build/

# Makefile
.PHONY: check build clean

check:
	./check.sh

build:
	python -m build --wheel

clean:
	rm -rf dist build src/*.egg-info

\end{lstlisting}

\begin{lstlisting}[language=bash]
# 红 → 绿
./check.sh                    # 改坏后：1 failed（红）；修复后：2 passed（绿）

# 构建与校验
make check
make build
sha256sum dist/greetlab_25100021005-0.1.0-py3-none-any.whl

# 内容聚焦提交（外部仓库根目录）
git add .gitignore
git commit -m "chore: ignore python cache and build artifacts"
git add week4/q/q16
git commit -m "week4 q16: scaffold greetlab from q13"
git add week4/q/q16/src/greetlab/cli.py
git commit -m "fix: greet with actual name instead of literal string"
git log --oneline


\end{lstlisting}

\begin{center}
{\large\bfseries 部分实验过程截图}
\end{center}

\resultimg{week\Week_check4.png}{题目四 实验过程截图:运行check.sh, ruff 照常通过，但pytest失败}
\resultimg{week\Week_check4_2.png}{题目四 实验过程截图:修复后通过makefile运行check和build，正常}
\resultimg{week\Week_check4_3.png}{题目四 实验过程截图:查看wheel的SHA-256值，并提交最终改动}


\begin{center}
{\large\bfseries 实验结果}
\end{center}

- 改坏后， ./check.sh 红；修复后绿（2 passed）

- 编写makefile，通过make build 产出 wheel，并显示 SHA-256 值

- 外部仓库三步提交（chore → scaffold → fix，仅 1 行改动）

\subsection*{课上阶段提交记录}
% 课上 4 题的 commit 记录截图（git log 或仓库 commits 列表）
\resultimgcap[0.9\textwidth]{week\Week_commit_kecheng.png}{课上阶段 commit 记录}

% ---------- 2. 课后练习与结果 ----------
\section{课后练习与结果}
% 说明：课后完成 >10 个实例，每个实例按「实验题目 → 代码 → 运行结果」组织。

\subsection*{实例 1}
\resetpic
% 实验题目：TODO 简述本实例的实验题目与要求
\textbf{实验题目：}此处填写实例一的实验题目与要求描述。

\begin{lstlisting}
# TODO: 实例1代码汇总
\end{lstlisting}

\textbf{运行结果：}

\resultimg{week\Week_ex1.png}{实例一 运行结果}

\subsection*{实例 2}
\resetpic
% 实验题目：TODO 简述本实例的实验题目与要求
\textbf{实验题目：}此处填写实例二的实验题目与要求描述。

\begin{lstlisting}
# TODO: 实例2代码汇总
\end{lstlisting}

\textbf{运行结果：}

\resultimg{week\Week_ex2.png}{实例二 运行结果}

\subsection*{实例 3}
\resetpic
% 实验题目：TODO 简述本实例的实验题目与要求
\textbf{实验题目：}此处填写实例三的实验题目与要求描述。

\begin{lstlisting}
# TODO: 实例3代码汇总
\end{lstlisting}

\textbf{运行结果：}

\resultimg{week\Week_ex3.png}{实例三 运行结果}

\subsection*{实例 4}
\resetpic
% 实验题目：TODO 简述本实例的实验题目与要求
\textbf{实验题目：}此处填写实例四的实验题目与要求描述。

\begin{lstlisting}
# TODO: 实例4代码汇总
\end{lstlisting}

\textbf{运行结果：}

\resultimg{week\Week_ex4.png}{实例四 运行结果}

\subsection*{实例 5}
\resetpic
% 实验题目：TODO 简述本实例的实验题目与要求
\textbf{实验题目：}此处填写实例五的实验题目与要求描述。

\begin{lstlisting}
# TODO: 实例5代码汇总
\end{lstlisting}

\textbf{运行结果：}

\resultimg{week\Week_ex5.png}{实例五 运行结果}

\subsection*{实例 6}
\resetpic
% 实验题目：TODO 简述本实例的实验题目与要求
\textbf{实验题目：}此处填写实例六的实验题目与要求描述。

\begin{lstlisting}
# TODO: 实例6代码汇总
\end{lstlisting}

\textbf{运行结果：}

\resultimg{week\Week_ex6.png}{实例六 运行结果}

\subsection*{实例 7}
\resetpic
% 实验题目：TODO 简述本实例的实验题目与要求
\textbf{实验题目：}此处填写实例七的实验题目与要求描述。

\begin{lstlisting}
# TODO: 实例7代码汇总
\end{lstlisting}

\textbf{运行结果：}

\resultimg{week\Week_ex7.png}{实例七 运行结果}

\subsection*{实例 8}
\resetpic
% 实验题目：TODO 简述本实例的实验题目与要求
\textbf{实验题目：}此处填写实例八的实验题目与要求描述。

\begin{lstlisting}
# TODO: 实例8代码汇总
\end{lstlisting}

\textbf{运行结果：}

\resultimg{week\Week_ex8.png}{实例八 运行结果}

\subsection*{实例 9}
\resetpic
% 实验题目：TODO 简述本实例的实验题目与要求
\textbf{实验题目：}此处填写实例九的实验题目与要求描述。

\begin{lstlisting}
# TODO: 实例9代码汇总
\end{lstlisting}

\textbf{运行结果：}

\resultimg{week\Week_ex9.png}{实例九 运行结果}

\subsection*{实例 10}
\resetpic
% 实验题目：TODO 简述本实例的实验题目与要求
\textbf{实验题目：}此处填写实例十的实验题目与要求描述。

\begin{lstlisting}
# TODO: 实例10代码汇总
\end{lstlisting}

\textbf{运行结果：}

\resultimg{week\Week_ex10.png}{实例十 运行结果}

\subsection*{实例 11}
\resetpic
% 实验题目：TODO 简述本实例的实验题目与要求
\textbf{实验题目：}此处填写实例十一的实验题目与要求描述。

\begin{lstlisting}
# TODO: 实例11代码汇总
\end{lstlisting}

\textbf{运行结果：}

\resultimg{week\Week_ex11.png}{实例十一 运行结果}

\subsection*{课后阶段提交记录}
% 课后练习的 commit 记录截图
\resultimgcap[0.9\textwidth]{week\Week_commit_kehou.png}{课后阶段 commit 记录}

% ---------- 3. 解题感悟 ----------
\section{解题感悟}
% TODO: 记录做题过程中的思考、遇到的问题与解决方式、收获与总结。
此处填写解题感悟。

\end{document}